\section{Source and Delivery Note}

This homily addresses an adult parish assembly for the Nineteenth Sunday
after Pentecost in the 1962 Roman Missal, general-calendar occurrence
4 October 2026. Its governing reading is ``The Invited People Clothed in
Charity'' in the expansive study, \work{The Feast and the New Person}.
The compatible account of truthful speech, restrained anger and generous
work comes from ``The Renewal of Hands, Words, and Desires.'' The concise
study, \work{Charity and the Renewed Life}, develops their comparison.
The connection across the propers is editorial synthesis; individual
patristic arguments retain their named authors.

The spoken body contains \textbf{1,301 words}, excluding the title and
all terminal apparatus. At 115--125 words per minute it takes approximately
10.4--11.3 minutes before additional pauses: an estimated 10--12-minute
delivery. A silent editorial read-through checked sense, transitions and
sentence length; no audible rehearsal or timed human delivery is claimed.
The examples are exhortations, not reports of personal experience.
The speech explains the Mass's prayers without supplying a new recited
prayer or a substitute liturgical translation. This is AI-assisted
preaching prose, not an ecclesiastical approval or an official liturgical text.

\section{References}
\begin{description}[style=nextline,leftmargin=0pt,labelsep=0pt,itemsep=.35em]
\item[Appointed Scripture and prayers]
\work{Missale Romanum}, Vatican editio typica, 1962, pp.~411--412,
nn.~1680--1681, 1684, 1686--1688: Collect, Eph 4:23--28,
Matt 22:1--14, Secret, Communion (Ps 118[119]:4--5) and
Postcommunion. \url{https://media.churchmusicassociation.org/pdf/missale62.pdf}.
The canonical study's collated texts retain the Pustet 1862 Latin
antecedent, pp.~351--352; historical English is Douay--Rheims--Challoner,
Gutenberg eBook 1581, and Cummiskey's \work{Roman Missal Translated
into the English Language for the Use of the Laity} (1861),
``Nineteenth Sunday after Pentecost,'' pp.~445--448.
The homily gives exposition, not an inventory of those texts.

\item[Augustine]
Sermon 90.4--6, 9--10, on the garment, growing charity and love of
enemies; NPNF I.6, English Sermon XL, Benedictine XC.
\work{Enarrationes in Psalmos} 118, opening Aleph exposition,
\S\S4--5, on command and petition; NPNF I.8, English heading Psalm 119.
\url{https://www.ccel.org/ccel/schaff/npnf106.html};
\url{https://www.ccel.org/ccel/schaff/npnf108.html}.

\item[Gregory the Great and John Chrysostom]
Gregory, \work{Homiliae in Evangelia} II.38.9--12, charity,
the Bridegroom and imperfect love; historical Latin transcription,
\url{https://la.wikisource.org/wiki/Homiliarum_in_Evangelia}.
Chrysostom, \work{Homilies on Ephesians} XIV, on 4:25--28,
truthful members, anger and labour for sharing; NPNF I.13,
\url{https://www.ccel.org/ccel/schaff/npnf113.html}.
Both are used in paraphrase.

\item[Blessed Ildefonso Schuster]
\work{The Sacramentary}, English edition, Burns Oates \& Washbourne,
1927, vol.~III, ``Nineteenth Sunday after Pentecost,''
pp.~171--174, especially the Collect, Secret and Postcommunion:
freedom for service, fruitful participation and healing grace.
Historical English inspected through the retained OCR; paraphrased.
\end{description}
